\chapter{L'art du soin}
\label{chap:l-art-du-soin}

La caravane était faite de permanents et de voyageurs de passage. Certains
restaient une journée, d'autres plusieurs semaines, parfois quelques mois.
Sans que Khaer Ghal et Kamatsu en aient encore pleinement conscience, chacun
d'eux pouvait laisser derrière lui quelque chose qui survivrait longtemps à
son départ.

Artis fut l'un de ceux-là.

Il voyageait avec la caravane depuis quelque temps déjà et les garçons avaient
rapidement compris qu'il savait soigner. Ils ne s'intéressèrent cependant
réellement à lui que lorsque Bartiméus, dont l'opinion de ses propres capacités
semblait très supérieure à la réalité, revint un jour dans un état inquiétant.
Nul ne sut exactement ce qu'il avait affronté. Au vu de ses blessures, quelque
chose de plus gros que lui avait probablement refusé de reconnaître sa
supériorité.

Kamatsu fut le premier à le trouver.

— Bartiméus !

Le chat roux avançait péniblement entre les chariots, le pelage sale et marqué
par plusieurs blessures. Lorsqu'il aperçut les deux garçons, il tenta encore de
conserver un semblant de dignité, mais renonça rapidement et se laissa
approcher.

— Qu'est-ce que t'as encore fait ? demanda Kamatsu en s'accroupissant près de
lui.

Bartiméus miaula faiblement. Khaer Ghal examina rapidement l'animal.

— Il faut trouver Artis.

Ils transportèrent immédiatement le chat jusqu'au guérisseur. Artis les fit
installer Bartiméus devant lui et commença par l'examiner. Il écarta
délicatement son pelage, observa plusieurs blessures, en toucha certaines avec
précaution et en évita d'autres. Kamatsu resta près de la tête du chat, une
main posée sur lui pour le rassurer.

Khaer Ghal, lui, commença progressivement à regarder autre chose.

Les mains d'Artis.

Le guérisseur nettoyait une plaie lorsqu'il s'interrompit et posa deux doigts
près de la blessure. Une faible lueur apparut autour de sa main, suffisamment
discrète pour ne pas éclairer toute la pièce, mais assez nette pour que Khaer
Ghal la remarque immédiatement. Il se pencha légèrement.

— C'était quoi ?

Artis ne releva pas les yeux.

— Un soin.

— Avec de la magie ?

— Oui.

Khaer Ghal fixa ses doigts.

— Comment tu fais ?

— Pour l'instant, je soigne le chat.

— Mais comment tu sais où toucher ?

Artis soupira doucement.

— Kerbal...

Khaer Ghal releva brièvement les yeux vers lui. Une nouvelle façon d'écorcher
son nom. Il n'en fit aucune remarque.

— Et pourquoi tu as commencé par celle-là ? Celle-ci est plus grande.

Il désigna une autre blessure.

— Plus grande ne veut pas dire plus grave.

Khaer Ghal fronça les sourcils et regarda de nouveau Bartiméus.

— Comment tu sais laquelle est la plus grave ?

Artis leva enfin les yeux vers lui.

— En apprenant à regarder.

Cette réponse eut exactement l'effet inverse de celui qu'il espérait peut-être.

— Regarder quoi ?

Kamatsu tourna la tête vers son ami.

— Laisse-le peut-être finir avant de lui poser toutes les questions.

Khaer Ghal se tut quelques secondes avant que sa curiosité ne reprenne le
dessus.

— Ça marche aussi sur les Orcs ?

Artis eut un sourire malgré lui.

— Oui.

— Et sur toutes les créatures ?

— Non.

— Pourquoi ?

Artis lui lança un regard.

— Kerbal...

— D'accord.

Cette fois, Khaer Ghal se tut réellement, mais continua d'observer chacun des
gestes du guérisseur.

%\illustration
%  {1200.png}
%  {L'art du soin}
%  {l-art-du-soin}

Lorsque Bartiméus fut hors de danger, Artis accepta finalement de répondre plus
longuement à ses questions. Il lui expliqua que soigner ne consistait pas
simplement à refermer une blessure : il fallait d'abord comprendre ce qui
s'était passé, reconnaître ce qui nécessitait une intervention immédiate et
savoir ce qu'il valait mieux ne pas toucher.

Il lui montra quelques gestes simples et lui expliqua comment nettoyer une
plaie, limiter un saignement ou reconnaître certains signes inquiétants. Khaer
Ghal écoutait avec la même attention qu'il accordait aux chasseurs lorsqu'ils
lui montraient une piste qu'il n'avait pas remarquée. Mais ce qui l'intriguait
le plus restait la magie.

— Je peux essayer ?

Artis le regarda.

— Si tu veux.

Khaer Ghal plaça sa main comme il l'avait vu faire et se concentra sur ses
doigts.

Rien ne se passa.

Il changea légèrement leur position.

Toujours rien.

— Ça ne marche pas.

— Pas aujourd'hui.

Khaer Ghal releva les yeux.

— Ça veut dire que ça peut marcher demain ?

Artis sourit.

— Ça veut surtout dire que tu devrais commencer par apprendre ce que tu peux
faire avec tes mains avant de t'inquiéter de ce que tu pourrais faire avec de
la magie.

La réponse ne satisfaisait qu'à moitié Khaer Ghal, mais il l'accepta. Dans les
jours qui suivirent, il revint régulièrement voir Artis. Le guérisseur lui
apprit quelques rudiments de soins, lui montra certaines plantes utiles et lui
expliqua surtout les erreurs à ne pas commettre. Khaer Ghal posa beaucoup de
questions, en oublia certaines réponses, en retint d'autres et essaya de
reproduire tout ce qu'on acceptait de lui montrer.

La magie, elle, refusa obstinément de se manifester.

Cela ne l'empêcha pas de continuer.

Lorsque Artis quitta finalement la caravane, Khaer Ghal était encore très loin
d'être un guérisseur. Il savait cependant nettoyer et protéger une blessure,
reconnaître quelques plantes utiles et, surtout, observer avant d'agir.

C'était peu.

Mais c'était davantage que ce qu'il savait avant de rencontrer Artis.